%% ===========================================================================
%%  Journal of Seasonality — submission (built on artaquest.cls).
%%  Reviewed by RUNNING it: the reviewer compiles this, fetches the open data,
%%  runs reproduce.py, and checks the headline numbers below reproduce.
%% ===========================================================================
\documentclass[11pt]{artaquest}

% Headline numbers — every one is printed by code/reproduce.py on the open data.
% They are kept as macros so the manuscript and the reproduction stay in lock-step.
\newcommand{\headlineN}{24}            % fields fitted
\newcommand{\headlineW}{1173}          % weeks
\newcommand{\medianR}{XX.X}            % median in-sample R^2 (%)
\newcommand{\meanR}{XX.X}              % mean in-sample R^2 (%)
\newcommand{\scaleEps}{X.X}            % max mean-R^2 spread across the global-scale sweep (pp)
\newcommand{\scaleFlat}{0.0X}          % mean-R^2 spread for s>=4 (flat-top regime, pp)
\newcommand{\pcaMean}{XX}              % data-optimal PCA ceiling, mean (%)
\newcommand{\baselineMean}{XX}         % no-astrology trend+annual baseline, mean (%)
\newcommand{\zodiacEps}{0.000X}        % R^2 spread across ayanamsa offsets

\title{Globally-Anchored, Robust Seasonality in Worldwide Search Interest, and the
       Scale-Invariance of Period-Sinc Kernels}
\runningtitle{Globally-anchored robust seasonality}

\author{Arash Ashrafnejad\affmark{1}}
\affiliation{\affmark{1}ArtaQuest Foundation}
\keywords{google trends; search interest; seasonality; robust regression; reproducibility; time series}

\dataavailability{The open dataset of \headlineN{} globally-anchored weekly search-interest series
  (\headlineW{} weeks, 2004--2026), together with the precomputed sidereal ephemeris, is openly
  available at \url{https://artaquest.org/wp-content/uploads/research/series2.json} and
  \url{https://artaquest.org/wp-content/uploads/research/ephemeris2.json}.}
\codeavailability{A self-contained reproduction bundle is openly available at
  \url{https://artaquest.org/papers/seasonality2/code.zip}; running \texttt{python3 reproduce.py}
  regenerates every headline number below (it needs only \texttt{numpy}, \texttt{scipy} and
  \texttt{scikit-learn}, and ships the data so no network is required).}

\begin{document}
\maketitle

\begin{abstract}
We revisit a transparent curve-fitting question --- does the \emph{timing} of worldwide curiosity, as
indexed by Google Trends, line up with the positions of celestial bodies? --- and harden it in three
ways, while keeping the same honest stance that correlation is not causation. First, we put every series
on a \emph{global} popularity scale: each query is measured in the same request against two fixed
reference words (``arta'' and ``quest''), so interest is comparable \emph{across} topics rather than each
being normalised to its own maximum. Second, we drop the cube-root variance-stabilising transform of
prior work and instead fit the \emph{raw} interest under a robust \emph{Huber} loss, moving the defence
against spikes from the data into the estimator. Third, we add a single global multiplier to the
period-scaled ``sinc'' kernel and ask whether tuning it helps. For $N=\headlineN{}$ single-word fields ---
a balanced design covering all twelve zodiac houses with both a noun and an adjective field each --- we
fit \headlineW{} weeks of stitched weekly interest (2004--2026) with the same fourteen-number model:
twelve sidereal bodies (the Moon taken synodically), an intercept, and one shared phase angle swept at
$1^\circ$. The median in-sample $R^2$ is $\medianR\%$ (mean $\meanR\%$). Our central methodological
finding is a \textbf{scale-invariance}: sweeping the global kernel multiplier over almost two orders of
magnitude moves the mean in-sample $R^2$ by only $\scaleEps$ percentage points, with no usable optimum, and
once every body is broad enough to enter the kernel's flat-top ($s\ge4$) the fit is flat to within
$\scaleFlat$ points. We give the reason analytically --- the slow outer bodies that carry the fit sit in
the kernel's quadratic flat-top, where rescaling the width merely rescales each feature's amplitude, which
the non-negative weights absorb \emph{exactly} --- so the single global scale is a redundant parameter, not
a useful one (the small residual variation is the fast inner bodies, which carry almost none of the fit). As before, three honest bounds (a data-optimal PCA ceiling, a no-astrology trend-plus-annual
baseline, and invariance to the zodiac convention) show the fit is a compact re-description of trend and
annual seasonality, not evidence of celestial influence.
\end{abstract}

\section{Introduction}
When millions of people type a word into a search box, the aggregate count is a noisy but honest record
of where collective attention drifts over time~\citep{choi2012predicting}. A companion study~\citep{ashrafnejad2026topics} fit two
decades of such series to a strict, physics-scaled model of sidereal cycles, and found --- after careful
bounding --- that the apparent celestial alignment is really trend and annual seasonality expressed in
orbital language. This paper does not relitigate that conclusion; it \emph{strengthens the method} and
reports one clean new result about the model itself.

Our contribution is threefold. (i)~\textbf{Global anchoring.} Google Trends normalises every query to its
own maximum, which destroys cross-topic comparability: a niche word and a household word both peak at
$100$. We instead measure every query in the \emph{same} three-term request against two fixed anchors,
recovering a single global interest scale (Section~\ref{sec:data}). (ii)~\textbf{Robust raw fitting.} We
abandon the cube-root transform and fit raw interest under a Huber loss~\citep{huber1964robust}, which
down-weights spikes without distorting the bulk of the data (Section~\ref{sec:model}).
(iii)~\textbf{A scale-invariance result.} We introduce a single global multiplier on the period-scaled
kernel and show --- analytically and empirically --- that the in-sample fit does not depend on it
(Section~\ref{sec:scale}). We keep the two self-disciplines that make the exercise hard to over-fit:
exactly fourteen numbers, and every feature constant a real orbital period.

\section{Data}
\label{sec:data}

\paragraph{A global popularity anchor.}
Trends returns each query rescaled so its own maximum is $100$, which makes absolute popularity
incomparable between queries. To recover a global scale we never request a word alone: every fetch is a
three-term comparison, \texttt{[query, ``arta'', ``quest'']}. The two anchors are deliberately chosen ---
``quest'' is a common, stable reference word, and ``arta'' a rarer one that provides a lower rung when a
query dwarfs ``quest''. We calibrate the fixed ratio between the two anchors once (median over a
multi-year overlap), cache it, and thereafter express every series as a percentage of the ``quest''
baseline. The result is one number per week that is comparable \emph{across} topics, not merely within
one. We report this global popularity alongside each field.

\paragraph{Stitching a 22-year weekly series.}
Trends returns \emph{weekly} resolution only for windows shorter than about five years, each renormalised
to its own maximum. To build one continuous weekly series over 2004--2026 we fetch many overlapping
$\sim$4-year weekly windows and stitch them onto the all-time monthly export as a backbone: rescale each
chunk to the backbone by least-squares through the origin, average the overlaps, then monthly-anchor each
month so its weekly mean matches the backbone while preserving within-month shape. We keep only series
whose pre-anchor stitch correlation is at least $0.95$. We then drop the most recent \num{52} weeks. This crop
is not arbitrary: it applies the fitting-power criterion of \citet{ashrafnejad2026recency}, who show that
Google Trends over-weights very recent activity and that in-sample fitting power climbs as the noisy
recent tail is removed, kneeing near one year. We confirmed the same knee on a held-out tuning sample
before fixing the crop at \num{52} weeks.

\paragraph{A balanced twelve-house design.}
Fields are single words drawn from the platform's topic catalogue, each tagged as a noun or an adjective.
We collect until the design is balanced: every one of the twelve zodiac houses (the model assigns each
field a best-fit sign) is represented by at least one noun field \emph{and} one adjective field --- the
$N=\headlineN{}$ fields analysed here. This guards against a sign distribution dominated by one
part-of-speech or one semantic cluster.

\section{Model}
\label{sec:model}

For a trial phase angle $\varphi$, swept in $1^\circ$ steps around the full $360^\circ$, each body $b$
contributes a single feature
\begin{equation}
  \mathrm{feature}_b(t) \;=\;
  \operatorname{sinc}\!\left(
    \frac{\,\Delta_\varphi\!\big(\lambda_b(t) - \varphi\big)\,}{s\,P_b}
  \right),
  \qquad
  \operatorname{sinc}(x) = \frac{\sin(\pi x)}{\pi x},
  \label{eq:feature}
\end{equation}
where $\lambda_b(t)$ is the body's sidereal longitude at time $t$ (Lahiri ayanamsa), $\Delta_\varphi(\cdot)$
is the wrapped angular difference in $(-180^\circ,180^\circ]$ expressed in radians, $P_b$ is the body's
orbital period in years, and $s$ is a \emph{single global scale} shared by all bodies (Section~\ref{sec:scale}).
Setting $s=1$ recovers the pure period-scaled kernel. We fit the raw (untransformed) weekly series as
\begin{equation}
  \mathrm{search}(t) \;\approx\; \beta_0 \;+\; \sum_{b=1}^{12} w_b \cdot \mathrm{feature}_b(t),
  \qquad w_b \ge 0,
  \label{eq:fit}
\end{equation}
with non-negative weights, by minimising a \textbf{Huber loss}~\citep{huber1964robust} rather than the
plain squared error. The Huber loss is quadratic for small residuals and linear beyond a knee set robustly
from the data (the scaled median absolute deviation), so occasional interest spikes are down-weighted
\emph{without} a non-linear transform of the data. This replaces the cube-root-plus-least-squares recipe
of prior work: robustness now lives in the estimator, and all $R^2$ values are reported on the raw
interest scale. The fourteen parameters are the twelve weights $w_b$, the intercept $\beta_0$, and the
shared phase $\varphi$; the angle of minimum Huber cost is the field's assigned \textbf{sign}. The Moon is
taken synodically (elongation from the Sun, period \SI{29.53}{\day}); the remaining bodies are the Sun,
Mercury, Venus, Mars, Jupiter, Saturn, Uranus, Neptune, Pluto, Chiron and the lunar node.

\section{The global scale is a redundant parameter}
\label{sec:scale}

The natural question once a global scale $s$ is in \eqref{eq:feature} is how to tune it. The answer,
perhaps surprisingly, is that one cannot: the in-sample fit is \emph{invariant} to $s$.

\paragraph{Empirically.}
Sweeping $s$ over $\{0.5, 1, 2, 4, 8, 16, 32\}$ and refitting every field under the model's non-negativity,
the mean in-sample $R^2$ moves by only $\scaleEps$ percentage points across the whole range, with no usable
optimum --- $s=1$ sits within noise of the best. The residual variation is confined to \emph{small} $s$,
where the fast inner bodies still carry oscillatory structure; for $s\ge4$, where every body has entered
the flat-top, the curve is flat to within $\scaleFlat$ percentage points. Tuning $s$ by in-sample $R^2$ ---
the right objective in this heavily underparametrised regime (fourteen numbers over $\sim$1{,}000 weekly
points, where overfitting is not a concern) --- therefore returns essentially a flat line, and we fix
$s=1$. (An \emph{unconstrained} fit would instead reward small $s$, by freely overfitting the fast bodies'
aliased high-frequency features; the non-negativity constraint, which suppresses those near-useless bodies,
is what makes the fit genuinely scale-free.)

\paragraph{Analytically.}
The reason is structural. The fit is carried by the slow outer bodies (Section~\ref{sec:results}), whose
angular travel over 22 years is small compared with the period-scaled kernel width $s P_b$. For these
bodies the sinc argument $x = \Delta_\varphi/(sP_b)$ is small, and
\begin{equation}
  \operatorname{sinc}(x) \;=\; 1 - \frac{(\pi x)^2}{6} + O(x^4)
  \;=\; 1 - \frac{\pi^2}{6}\,\frac{\Delta_\varphi^2}{s^2 P_b^2} + O(s^{-4}).
  \label{eq:flat}
\end{equation}
Changing $s$ in \eqref{eq:flat} rescales the feature's \emph{deviation} from $1$ by the constant factor
$s^{-2}$ but leaves its \emph{shape} --- proportional to $\Delta_\varphi^2(t)$ --- unchanged. A linear
model with a free intercept and a free (here, non-negative) weight absorbs that constant exactly: the
optimum simply rescales $w_b \mapsto s^2 w_b$, leaving the fitted values, the residuals and hence $R^2$
unchanged. In other words, in the kernel's quadratic flat-top --- where the slow bodies that dominate the
fit all live --- a global width scale is degenerate with the weights. Only fast bodies, whose features
have genuine oscillatory structure within the window, can break the degeneracy, and they contribute too
little to move the aggregate fit. This is a clean negative result: a single global kernel scale is a
redundant knob, and the parameter-free period scaling ($s=1$) is the right default.

\section{Results}
\label{sec:results}

On the raw, non-negative, Huber-robust model the in-sample fit is strong: the \textbf{median $R^2$ is
$\medianR\%$} and the \textbf{mean $\meanR\%$} across the $N=\headlineN{}$ fields. A representative fit and
the distribution of best-fit signs are shown in \cref{fig:fit,fig:signs}.

\IfFileExists{seasonality2_fig1.png}{%
  \begin{figure}[ht]\centering
    \includegraphics[width=0.80\textwidth]{seasonality2_fig1.png}
    \caption{A representative field: the stitched, globally-anchored weekly interest (observed) and the
      fourteen-parameter Huber fit. The explained variation is mostly slow multi-year trend plus the
      annual rhythm.}
    \label{fig:fit}
  \end{figure}}{}

\paragraph{Where the fit comes from.}
A leave-one-flat variance decomposition (flatten each body's feature to its mean, refit, record the
non-negative $R^2$ drop, normalise to sum to the full $R^2$) gives the same picture as prior work: the fit
is carried by the \emph{slow outer bodies} (Saturn, Uranus, Neptune, Pluto and the node), which together
form a smooth multi-year \emph{trend} basis, with the Sun supplying the once-a-year \emph{seasonality}.
The fast inner bodies contribute little --- which is exactly why the global scale of Section~\ref{sec:scale}
is degenerate.

\IfFileExists{seasonality2_fig2.png}{%
  \begin{figure}[ht]\centering
    \includegraphics[width=0.70\textwidth]{seasonality2_fig2.png}
    \caption{Distribution of best-fit signs across the $N=\headlineN{}$ fields (the swept phase angle of
      minimum Huber cost). The balanced design ensures every house carries both a noun and an adjective
      field.}
    \label{fig:signs}
  \end{figure}}{}

\section{What the fit is, and is not (honest bounds)}
\label{sec:bounds}
We bound the fit from three directions, exactly as a sceptic should.

\paragraph{(a) A data-optimal ceiling.}
Principal-component analysis gives the best possible fit of a given size. The best
\emph{thirteen}-coefficient model (twelve principal components plus a per-field intercept) on the raw
weekly data explains about $\pcaMean\%$ of the variance (released code). Our fourteen-number model sits
below this ceiling: it leaves real structure unexplained and is, if anything, weaker than a size-matched
optimal basis --- consistent with a coarse trend-and-seasonal re-description rather than any
astrology-specific power.

\paragraph{(b) A plain baseline does at least as well.}
A no-astrology baseline --- a smooth polynomial trend plus an annual harmonic, with comparable degrees of
freedom and \emph{zero} celestial input --- explains about $\baselineMean\%$ (mean) of the variance, at
least as much as the full model. The explanatory work is trend and seasonality, not celestial specificity.

\paragraph{(c) The zodiac system does not matter.}
Across a sweep of ayanamsa offsets (the choice of sidereal convention) the mean $R^2$ varies by only
$\zodiacEps$. Switching convention shifts every longitude by a near-constant offset, which the single
shared phase $\varphi$ absorbs; the synodic Moon is offset-invariant to begin with.

\paragraph{Conclusion.}
Bounds (a)--(c) repeat the companion study's honest reading under the new, more robust estimator: the
sidereal alignment is a compact re-description of trend and annual seasonality in orbital language, not
evidence of celestial influence on human curiosity.

\section{Limitations}
\begin{itemize}
  \item \textbf{Descriptive, not predictive.} Every $R^2$ is in-sample; we make no forecasting claim.
  \item \textbf{Alignment is not causation.} Section~\ref{sec:bounds} shows the fit is plain trend and
        seasonality with no astrology.
  \item \textbf{A modest sample, balanced by design.} $N=\headlineN{}$ fields are enough to balance the
        twelve-house $\times$ \{noun, adjective\} design but remain a small, single-language (English)
        sample; the headline $R^2$ should be read as descriptive of this set, not a population estimate.
  \item \textbf{Reproducibility as the safeguard.} Because the claims are easy to over-read, we publish
        the globally-anchored data and the code so any reader can recompute every number.
\end{itemize}

\section*{Reproducibility}
The accompanying script (\texttt{reproduce.py}) regenerates, from the open data, the headline median and
mean $R^2$, the sign distribution, the global-scale invariance sweep of Section~\ref{sec:scale}, the
per-body decomposition, and the three bounds of Section~\ref{sec:bounds}.

\declarations
\authorcontributions{A.A.\ designed the analysis, wrote the code and the manuscript, and verified the
  reproduction.}
\competinginterests{The author declares no competing interests.}
\funding{This work received no external funding.}
\acknowledgements{This work uses Google Trends data, the Swiss Ephemeris, and the open-source scientific
  Python stack~\citep{harris2020numpy,virtanen2020scipy,pedregosa2011scikit}. The author thanks the
  ArtaScience reviewer for running the code and checking that the numbers reproduce.}

\bibliography{references}
\end{document}
